% Generated by scripts/assemble.py; corrections belong in source/corrections.json.
\examhead{96}{南京理工大学 2024 年数学分析试卷}{95}
% p096-b20; source PDF page 96
\par\Needspace{5\baselineskip}
\blockanchor{p096-b20}
\majorquestion{一}(15 分) 解答如下问题:\par
% p096-b21; source PDF page 96
\blockanchor{p096-b21}
(1) 求极限  $\lim_{x\to\infty}\left(\cos\frac{1}{x}\right)^{x^{2}-1}$ .\par
% p096-b22; source PDF page 96
\blockanchor{p096-b22}
(2) 设 $z = x^{2} + y^{2}$, 其中 $y = f(x)$ 是由曲线 $x^{2} - xy + y^{2} = 1$ 确定的隐函数, 求 $\frac{\mathrm{d}z}{\mathrm{d}x}$.\par
% p096-b23; source PDF page 96
\par\Needspace{5\baselineskip}
\blockanchor{p096-b23}
\majorquestion{二}(15 分) 解答如下问题:\par
% p096-b24; source PDF page 96
\blockanchor{p096-b24}
(1) 判断 $\sum_{n=1}^{\infty} \frac{1}{\left[\ln \left(\tan \frac{1}{n}\right)\right]^3}$ 的敛散性.\par
% p096-b25; source PDF page 96
\blockanchor{p096-b25}
(2) 判断  $\int_{1}^{+\infty}\frac{\sin x^{2}}{1+x^{p}}\mathrm{d}x\ (p\geq0)$  的敛散性.\par
% p096-b26; source PDF page 96
\par\Needspace{5\baselineskip}
\blockanchor{p096-b26}
\majorquestion{三}(15 分) 设 $f(x)$ 在 $\mathbb{R}$ 上二阶可导, 且 $f'(0) = 0$, 证明: 存在 $\xi \in (0,1)$, 使得 $f'(\xi) = (\xi - 1)^2 f''(\xi)$.\par
% p097-b00; source PDF page 97
\par\Needspace{5\baselineskip}
\blockanchor{p097-b00}
\majorquestion{四}(15 分) 设 $f(x) = \sum_{n=1}^{\infty} ne^{-nx}$.\par
% p097-b01; source PDF page 97
\blockanchor{p097-b01}
(1) 讨论 $f(x)$ 的连续区间.\par
% p097-b02; source PDF page 97
\blockanchor{p097-b02}
(2) 计算 $\int_{1}^{2} f(x) \, \mathrm{d}x$.\par
% p097-b03; source PDF page 97
\par\Needspace{5\baselineskip}
\blockanchor{p097-b03}
\majorquestion{五}(15 分) 设 $f(x)$ 在 $\mathbb{R}$ 上连续, 且 $\lim_{x \to \infty} [f(x) - \lambda x] = 0$, 其中 $\lambda$ 为常数, 证明: $f(x)$ 在 $\mathbb{R}$ 上一致连续.
\majorquestion{六}(15 分) 解答如下问题:\par
% p097-b04; source PDF page 97
\blockanchor{p097-b04}
(1) 求 $f(x) = \arctan \frac{x + 3}{x - 3}$ 的麦克劳林级数展开式, 并讨论收敛区间.\par
% p097-b05; source PDF page 97
\blockanchor{p097-b05}
(2) 求 $f(x) = x^3$ 在 $[0, \pi]$ 上的余弦级数.\par
% p097-b06; source PDF page 97
\par\Needspace{5\baselineskip}
\par\Needspace{6\baselineskip}
\blockanchor{p097-b06}
\majorquestion{七}(15 分) 计算三重积分\par
\nopagebreak[4]
% p097-b07; source PDF page 97
\blockanchor{p097-b07}
\displayformula{\iiint_ {V} z ^ {2} \mathrm{d} x \mathrm{d} y \mathrm{d} z.}
% p097-b08; source PDF page 97
\blockanchor{p097-b08}
其中 $V: x^{2} + y^{2} + z^{2} \leq a^{2}, x^{2} + y^{2} \leq ax, a > 0.$\par
% p097-b09; source PDF page 97
\par\Needspace{5\baselineskip}
\par\Needspace{6\baselineskip}
\blockanchor{p097-b09}
\majorquestion{八}(15 分) 计算曲面积分\par
\nopagebreak[4]
% p097-b10; source PDF page 97
\blockanchor{p097-b10}
\displayformula{I = \iint_ {S} x ^ {2} \mathrm{d} y \mathrm{d} z + y ^ {2} \mathrm{d} z \mathrm{d} x + z ^ {2} \mathrm{d} x \mathrm{d} y.}
% p097-b11; source PDF page 97
\blockanchor{p097-b11}
其中 $S$ 为 $x^{2} + y^{2} = z^{2}, 0 \leq z \leq 1$, 取曲面下侧.\par
% p097-b12; source PDF page 97
\par\Needspace{5\baselineskip}
\par\Needspace{6\baselineskip}
\blockanchor{p097-b12}
\majorquestion{九}(15 分) 计算曲线积分\par
\nopagebreak[4]
% p097-b13; source PDF page 97
\blockanchor{p097-b13}
\displayformula{\lim _ {d (\Omega) \rightarrow 0} \frac {1}{S (\Omega)} \oint_ {L} (\boldsymbol {F}, \boldsymbol {n}) d s.}
% p097-b14; source PDF page 97
\blockanchor{p097-b14}
其中 $\Omega$ 是包含原点且由简单封闭的光滑曲线 $C$ 围成的区域, $d(\Omega)$ 为 $\Omega$ 的直径, $S(\Omega)$ 为 $\Omega$ 的面积, $\pmb{F} = (P(x,y), Q(x,y))$ 为区域 $\Omega + C$ 上连续可微的向量函数, $\pmb{n}$ 为曲线 $C$ 的单位外法向量.\par
% p097-b15; source PDF page 97
\par\Needspace{5\baselineskip}
\blockanchor{p097-b15}
\majorquestion{十}(15 分) 用有限覆盖定理证明致密性定理, 即任意有界数列 $\{x_{n}\}$ 必定有收敛子列.\par
